\SourceLeaf{239}%
and bear in mind that the resulting values of the differentials \(\mathrm{d}t,\mathrm{d}u\) must be equal, independently of the differentials \(\mathrm{d}x,\mathrm{d}y\), to the quantities \(T\mathrm{d}x+U\mathrm{d}y\) and \(U\mathrm{d}x+V\mathrm{d}y\), respectively, we will find, after some fairly evident transformations:
\[
\begin{aligned}
C^3T&=\alpha Ab'b'+\beta Bb'b'+\gamma Cb'b'\\
&\quad-2\alpha'Abb'-2\beta'Bbb'-2\gamma'Cbb'\\
&\quad+\alpha''Abb+\beta''Bbb+\gamma''Cbb\\
C^3U&=-\alpha Aa'b'-\beta Ba'b'-\gamma Ca'b'\\
&\quad+\alpha'A(ab'+ba')+\beta'B(ab'+ba')+\gamma'C(ab'+ba')\\
&\quad-\alpha''Aab-\beta''Bab-\gamma''Cab\\
C^3V&=\alpha Aa'a'+\beta Ba'a'+\gamma Ca'a'\\
&\quad-2\alpha'Aaa'-2\beta'Baa'-2\gamma'Caa'\\
&\quad+\alpha''Aaa+\beta''Baa+\gamma''Caa
\end{aligned}
\]
If, then, for brevity we put
\[
\begin{aligned}
A\alpha+B\beta+C\gamma&=D\qquad .\ .\ .\ .\ .\ .\ .\ .\ .\ .\quad(1)\\
A\alpha'+B\beta'+C\gamma'&=D'\qquad .\ .\ .\ .\ .\ .\ .\ .\ .\ .\quad(2)\\
A\alpha''+B\beta''+C\gamma''&=D''\qquad .\ .\ .\ .\ .\ .\ .\ .\ .\ .\quad(3)
\end{aligned}
\]
we have
\[
\begin{aligned}
C^3T&=Db'b'-2D'bb'+D''bb\\
C^3U&=-Da'b'+D'(ab'+ba')-D''ab\\
C^3V&=Da'a'-2D'aa'+D''aa
\end{aligned}
\]
Expanding this, we find
\[
C^6(TV-UU)=(DD''-D'D')(ab'-ba')^2=(DD''-D'D')CC
\]
and hence the formula for the measure of curvature
\[
k=\frac{DD''-D'D'}{(AA+BB+CC)^2}
\]
\Article{11}
We shall now build on the formula just obtained to derive another, which should rank among the most fruitful theorems in the theory of curved surfaces. Let us introduce the following notation:
